% !TeX program = xelatex
\documentclass[11pt,a4paper]{article}

\usepackage{fontspec}
\setmainfont{Times New Roman}

\usepackage[a4paper,margin=2.5cm]{geometry}
\usepackage{array}
\usepackage{longtable}

\renewcommand{\tablename}{Tabela}

\setlength{\parindent}{0pt}
\setlength{\parskip}{4pt}

\title{\vspace{-2em}{\Large\bfseries UML --- Unified Modeling Language}}
\author{}
\date{}

\begin{document}
\maketitle
\vspace{-3em}

\section*{O que é}

UML (Unified Modeling Language, ou Linguagem de Modelagem Unificada) é uma
linguagem visual padronizada para especificar, visualizar, construir e
documentar os artefatos de um sistema de software. Não é uma linguagem de
programação: não executa, não compila --- serve para representar
graficamente a estrutura e o comportamento de um sistema antes, durante e
depois da implementação.

Mantida pela OMG (Object Management Group), surgiu em 1997 da unificação de
três notações concorrentes (Booch, OMT de Rumbaugh, e OOSE de Jacobson), os
chamados ``três amigos'' da modelagem orientada a objetos.

\section*{Para que serve}

\begin{itemize}
  \item \textbf{Comunicação}: uma notação comum entre desenvolvedores,
    arquitetos, stakeholders e professores/orientadores --- todos leem o
    mesmo desenho.
  \item \textbf{Documentação}: registra decisões de projeto (estrutura de
    classes, fluxo de casos de uso, arquitetura de implantação) de forma
    durável, independente do código-fonte.
  \item \textbf{Análise e projeto}: ajuda a pensar a solução antes de
    codificar --- identificar entidades, responsabilidades,
    relacionamentos e fluxos.
  \item \textbf{Rastreabilidade}: liga requisitos a elementos de projeto,
    facilitando auditoria e manutenção.
\end{itemize}

\section*{As duas grandes famílias de diagramas}

UML define 14 tipos de diagrama, divididos em dois grupos.

\subsection*{1. Diagramas Estruturais (o que o sistema \emph{é})}

Mostram a organização estática dos elementos.

\begin{longtable}{p{4.2cm}p{10.3cm}}
  \hline
  \textbf{Diagrama} & \textbf{Representa} \\
  \hline
  \endhead
  \hline
  \endfoot
  Classes & Classes, atributos, métodos e relacionamentos (associação, herança, agregação, composição) \\
  Objetos & Instâncias concretas de classes em um momento específico \\
  Componentes & Módulos/componentes de software e suas interfaces \\
  Implantação (Deployment) & Distribuição física do software em nós de hardware/infraestrutura \\
  Pacotes & Agrupamento e dependência entre pacotes/namespaces \\
  Estrutura Composta & Estrutura interna de uma classe/componente e suas colaborações \\
  Perfil & Extensões e estereótipos customizados da própria UML \\
\end{longtable}

\subsection*{2. Diagramas Comportamentais (o que o sistema \emph{faz})}

Mostram interação, fluxo e mudança de estado ao longo do tempo.

\begin{longtable}{p{4.2cm}p{10.3cm}}
  \hline
  \textbf{Diagrama} & \textbf{Representa} \\
  \hline
  \endhead
  \hline
  \endfoot
  Casos de Uso & Funcionalidades do sistema sob a ótica de atores externos \\
  Sequência & Ordem temporal de troca de mensagens entre objetos \\
  Atividades & Fluxo de controle/processo, semelhante a um fluxograma \\
  Máquina de Estados & Estados possíveis de um objeto e as transições entre eles \\
  Comunicação & Interações entre objetos com foco na topologia (quem fala com quem) \\
  Visão Geral de Interação & Combina diagrama de atividades com diagramas de sequência \\
  Tempo (Timing) & Comportamento de objetos ao longo de uma linha do tempo \\
\end{longtable}

\section*{Elementos e notações essenciais}

\begin{itemize}
  \item \textbf{Classe}: retângulo dividido em três compartimentos ---
    nome, atributos, métodos. Visibilidade indicada por \texttt{+}
    (público), \texttt{-} (privado), \texttt{\#} (protegido), \texttt{\~{}}
    (pacote).
  \item \textbf{Relacionamentos}:
  \begin{itemize}
    \item \emph{Associação} (linha simples) --- uma classe usa/conhece a outra.
    \item \emph{Agregação} (losango vazio) --- relação ``todo-parte'' fraca;
      a parte existe independente do todo.
    \item \emph{Composição} (losango cheio) --- relação ``todo-parte''
      forte; a parte não existe sem o todo.
    \item \emph{Herança/Generalização} (seta com ponta vazia) ---
      especialização de uma classe base.
    \item \emph{Realização/Implementação} (seta tracejada com ponta vazia)
      --- classe implementa uma interface.
    \item \emph{Dependência} (seta tracejada) --- uma classe depende de
      outra, sem associação estrutural fixa.
  \end{itemize}
  \item \textbf{Multiplicidade}: indica quantas instâncias participam da
    relação (\texttt{1}, \texttt{0..1}, \texttt{1..*}, \texttt{0..*}, etc.).
  \item \textbf{Ator}: boneco (\emph{stick figure}) em diagramas de caso de
    uso, representa um usuário ou sistema externo que interage com o
    sistema.
  \item \textbf{Nó (node)}: cubo em diagramas de implantação, representa um
    recurso computacional físico ou virtual (servidor, dispositivo,
    container).
\end{itemize}

\section*{Por que usar em um TCC / projeto de software}

\begin{itemize}
  \item Formaliza requisitos funcionais em \textbf{casos de uso}.
  \item Formaliza a arquitetura estática em \textbf{diagrama de classes}.
  \item Formaliza a arquitetura física/implantação em \textbf{diagrama de
    implantação}.
  \item Dá rastreabilidade entre requisito $\rightarrow$ classe/método
    $\rightarrow$ execução, essencial para bancas e revisões técnicas.
  \item É a notação mais reconhecida academicamente para esse tipo de
    documentação, o que facilita avaliação por orientadores.
\end{itemize}

\section*{Ferramentas comuns}

PlantUML, Mermaid, draw.io/diagrams.net, StarUML, Visual Paradigm,
Enterprise Architect --- todas geram os mesmos diagramas UML a partir de
texto (PlantUML/Mermaid) ou de edição visual (draw.io, StarUML).

\section*{Referência}

OMG UML Specification --- \texttt{https://www.omg.org/spec/UML/}

\end{document}
